% GENERATED FILE. Edit canonical lessons, then run npm run generate:books.
% canonical-source-hash: fnv1a64:246b4be36aa5f04b
% canonical-lessons: ES-C486-afeitarse, ES-C486-toser, ES-C486-respirar, ES-C486-aspirina, ES-C486-vacuna

\chapter{En la consulta}
\label{ch:es-en-la-consulta}
\hlchaptermodality{486}

\begin{chapteropening}
\textbf{By the end of this chapter:} \emph{I can follow a doctor's advice about coughing, breathing, an aspirin and a vaccine, and read a note about shaving before an operation.}

Complete the last lesson of chapter 486: I can follow a doctor's advice about coughing, breathing, an aspirin and a vaccine, and read a note about shaving before an operation.
\end{chapteropening}

\section[afeitarse]{afeitarse — to shave}

\label{lesson:ES-C486-afeitarse}



\begin{quote}
Before the new one: say the Spanish for to dry, then the Spanish for the rubbish.
\end{quote}

\subsection*{You'll want to know: afeitarse}
\textbf{afeitarse} — \textquotedblleft{}to shave\textquotedblright{}.

\begin{quote}
Mi padre se afeita cada mañana.
\end{quote}

\begin{quote}
\emph{My father shaves every morning.}
\end{quote}

Reflexive, because you shave yourself.

\begin{cousinweb}
\textbf{afeitar} $\leftarrow$ Latin \emph{affectare}, \textquotedblleft{}to strive after\textquotedblright{} (Lewis \& Short). In Spanish it first meant \textquotedblleft{}to adorn, to make up\textquotedblright{}, and only later \textquotedblleft{}to shave\textquotedblright{}.
\end{cousinweb}

\begin{grammarlens}[title={What you've built}]
One more word for reading notices, forms and messages.
\end{grammarlens}

\subsection*{Your turn}
\begin{itemize}
\raggedright
  \item \emph{Say it:} \emph{afeitarse}
  \item \emph{Say it:} \emph{afeitarse}, once more
  \item \emph{Say it:} a sentence of your own with \emph{afeitarse}
  \item \emph{Recall:} say the Spanish for to dry, then the Spanish for the rubbish, then say \emph{afeitarse} again
\end{itemize}

\begin{tcolorbox}[breakable,colback=teal!4,colframe=teal!35!black,title={Before you move on}]
What does afeitarse mean? (\textquotedblleft{}To shave\textquotedblright{}.) Say it once more.
\end{tcolorbox}
\section[toser]{toser — to cough}

\label{lesson:ES-C486-toser}



\begin{quote}
Before the new one: say the Spanish for the rubbish, then the Spanish for to shave.
\end{quote}

\subsection*{You'll want to know: toser}
\textbf{toser} — \textquotedblleft{}to cough\textquotedblright{}.

\begin{quote}
El niño tose mucho por la noche.
\end{quote}

\begin{quote}
\emph{The child coughs a lot at night.}
\end{quote}

\emph{La tos} is the cough.

\begin{cousinweb}
\textbf{toser} $\leftarrow$ Latin \emph{tussire}, \textquotedblleft{}to cough\textquotedblright{} (EtymDB 2.1) — inherited: the \emph{u} became \emph{o}. \emph{La tos}, \textquotedblleft{}the cough\textquotedblright{}, is Latin \emph{tussis}.
\end{cousinweb}

\begin{grammarlens}[title={What you've built}]
One more word for reading notices, forms and messages.
\end{grammarlens}

\subsection*{Your turn}
\begin{itemize}
\raggedright
  \item \emph{Say it:} \emph{toser}
  \item \emph{Say it:} \emph{toser}, once more
  \item \emph{Say it:} a sentence of your own with \emph{toser}
  \item \emph{Recall:} say the Spanish for the rubbish, then the Spanish for to shave, then say \emph{toser} again
\end{itemize}

\begin{tcolorbox}[breakable,colback=teal!4,colframe=teal!35!black,title={Before you move on}]
What does toser mean? (\textquotedblleft{}To cough\textquotedblright{}.) Say it once more.
\end{tcolorbox}
\section[respirar]{respirar — to breathe}

\label{lesson:ES-C486-respirar}



\begin{quote}
Before the new one: say the Spanish for to shave, then the Spanish for to cough.
\end{quote}

\subsection*{You'll want to know: respirar}
\textbf{respirar} — \textquotedblleft{}to breathe\textquotedblright{}.

\begin{quote}
Respire hondo, por favor.
\end{quote}

\begin{quote}
\emph{Breathe deeply, please.}
\end{quote}

\begin{cousinweb}
\textbf{respirar} $\leftarrow$ Latin \emph{respirare}, \textquotedblleft{}to breathe out, to breathe again\textquotedblright{} (EtymDB 2.1): \emph{re-} and \emph{spirare}, \textquotedblleft{}to breathe\textquotedblright{}.
\end{cousinweb}

\begin{grammarlens}[title={What you've built}]
One more word for reading notices, forms and messages.
\end{grammarlens}

\subsection*{Your turn}
\begin{itemize}
\raggedright
  \item \emph{Say it:} \emph{respirar}
  \item \emph{Say it:} \emph{respirar}, once more
  \item \emph{Say it:} a sentence of your own with \emph{respirar}
  \item \emph{Recall:} say the Spanish for to shave, then the Spanish for to cough, then say \emph{respirar} again
\end{itemize}

\begin{tcolorbox}[breakable,colback=teal!4,colframe=teal!35!black,title={Before you move on}]
What does respirar mean? (\textquotedblleft{}To breathe\textquotedblright{}.) Say it once more.
\end{tcolorbox}
\section[la aspirina]{la aspirina — the aspirin}

\label{lesson:ES-C486-aspirina}



\begin{quote}
Before the new one: say the Spanish for to cough, then the Spanish for to breathe.
\end{quote}

\subsection*{You'll want to know: la aspirina}
\textbf{la aspirina} — \textquotedblleft{}the aspirin\textquotedblright{}.

\begin{quote}
Tómate una aspirina con agua.
\end{quote}

\begin{quote}
\emph{Take an aspirin with water.}
\end{quote}

\begin{cousinweb}
\textbf{aspirina} $\leftarrow$ German \emph{Aspirin}, a brand name coined by the Bayer company in 1899: \emph{a-} for acetyl, \emph{-spir-} for \emph{Spiraea}, the meadowsweet plant, and the chemists' ending \emph{-in}.
\end{cousinweb}

\begin{grammarlens}[title={What you've built}]
One more word for reading notices, forms and messages.
\end{grammarlens}

\subsection*{Your turn}
\begin{itemize}
\raggedright
  \item \emph{Say it:} \emph{la aspirina}
  \item \emph{Say it:} \emph{la aspirina}, once more
  \item \emph{Say it:} a sentence of your own with \emph{la aspirina}
  \item \emph{Recall:} say the Spanish for to cough, then the Spanish for to breathe, then say \emph{la aspirina} again
\end{itemize}

\begin{tcolorbox}[breakable,colback=teal!4,colframe=teal!35!black,title={Before you move on}]
What does la aspirina mean? (\textquotedblleft{}The aspirin\textquotedblright{}.) Say it once more.
\end{tcolorbox}
\section[la vacuna]{la vacuna — the vaccine}

\label{lesson:ES-C486-vacuna}



\begin{quote}
Before the new one: say the Spanish for to breathe, then the Spanish for the aspirin.
\end{quote}

\subsection*{You'll want to know: la vacuna}
\textbf{la vacuna} — \textquotedblleft{}the vaccine\textquotedblright{}.

\begin{quote}
La vacuna es gratis para los mayores.
\end{quote}

\begin{quote}
\emph{The vaccine is free for older people.}
\end{quote}

\emph{Vacunarse} — to get vaccinated.

\begin{cousinweb}
\textbf{vacuna} is built on \emph{vacuno}, \textquotedblleft{}of cattle\textquotedblright{} (EtymDB 2.1), from Latin \emph{vacca}, \textquotedblleft{}cow\textquotedblright{}: the first vaccine, in 1796, was made from cowpox.
\end{cousinweb}

\begin{grammarlens}[title={What you've built}]
One more word for reading notices, forms and messages.
\end{grammarlens}

\subsection*{Your turn}
\begin{itemize}
\raggedright
  \item \emph{Say it:} \emph{la vacuna}
  \item \emph{Say it:} \emph{la vacuna}, once more
  \item \emph{Say it:} a sentence of your own with \emph{la vacuna}
  \item \emph{Recall:} say the Spanish for to breathe, then the Spanish for the aspirin, then say \emph{la vacuna} again
\end{itemize}

\begin{tcolorbox}[breakable,colback=teal!4,colframe=teal!35!black,title={Before you move on}]
What does la vacuna mean? (\textquotedblleft{}The vaccine\textquotedblright{}.) Say it once more.
\end{tcolorbox}
